\chapter{FracMinHash Sketching}
\label{ch:sketching}

This chapter specifies \code{FracMinHash::sketch}, the single most bit-sensitive routine in the
whole system: every downstream index, seed, and score depends on hashes being computed exactly
this way, since reference and read sketches must use an \emph{identical} hash function to be
comparable at all.

\section{Inputs and outputs}

\begin{center}
\begin{tabular}{@{}p{2.6cm}p{4.5cm}p{7cm}@{}}
\toprule
& \textbf{Type} & \textbf{Meaning} \\
\midrule
\textit{Input} \code{s} & raw ASCII nucleotide bytes (upper- or lower-case \code{A/C/G/T};
  anything else, see \S\ref{sec:sketch-nonacgt}) & The sequence to sketch: one reference segment,
  or one read. \\
\textit{Config} \code{k} & positive integer & $k$-mer length (CLI \code{-k}, default 15). \\
\textit{Config} \code{h\_frac} & $r \in (0,1]$ & Keep-fraction of the hash space (CLI \code{-r},
  default 0.05). \\
\textit{Output} & \code{Vec<Kmer>} & Every $k$-mer of \code{s} whose canonical hash clears the
  threshold, in increasing position order, each tagged with its hash and winning strand
  (\S\ref{sec:type-kmer}). Empty if $|s| < k$. \\
\bottomrule
\end{tabular}
\end{center}

\section{Lookup tables: per-nucleotide hash contributions}
\label{sec:sketch-lut}

The rolling hash is built from two 256-entry lookup tables, \code{LUT\_fw} and \code{LUT\_rc},
indexed by raw ASCII byte value, giving each of the four nucleotides (in both upper and lower
case) a fixed 64-bit contribution:

\begin{center}
\begin{tabularx}{\textwidth}{@{}llX@{}}
\toprule
Base & \code{LUT\_fw} contribution (hex) & \code{LUT\_rc} contribution \\
\midrule
A/a & \code{0x3c8bfbb395c60474} & = \code{LUT\_fw['T']} \\
C/c & \code{0x3193c18562a02b4c} & = \code{LUT\_fw['G']} \\
G/g & \code{0x20323ed082572324} & = \code{LUT\_fw['C']} \\
T/t & \code{0x295549f54be24456} & = \code{LUT\_fw['A']} \\
\bottomrule
\end{tabularx}
\end{center}

That is, a base's reverse-complement contribution is defined as its \emph{complementary} base's
forward contribution --- exactly encoding the biological complement rule (A$\leftrightarrow$T,
C$\leftrightarrow$G) at the hash-contribution level, so that no explicit ``complement a
character'' step is ever needed; complementation falls out of table lookup alone. All other byte
values (N, ambiguity codes, whitespace, etc.) have \emph{no defined contribution} in the reference
source.

\begin{bugbox}
The reference C++ declares \code{LUT\_fw}/\code{LUT\_rc} as plain fixed-size arrays and
initializes only the 8 entries above; the remaining 248 entries are \textbf{uninitialized stack
memory} at the time \code{sketch()} first reads them. For any input byte other than
\texttt{ACGTacgt} (most commonly \texttt{N}), the hash computation silently incorporates garbage.
This is undefined behavior, not merely ``an unspecified but fixed value'' --- the actual bits
read depend on whatever happened to be on the stack. \textbf{Recommended fix for a from-scratch
implementation:} zero-initialize both tables before filling in the 8 real entries, so that any
non-ACGT byte deterministically contributes a hash of $0$. This changes no behavior for any
ACGT-only input (which is the overwhelming majority of real sequence data) and removes the UB.
\end{bugbox}

\section{The position semantics of \code{Kmer.r}}
\label{sec:sketch-r-semantics}

Before giving the recurrence, it is worth nailing down precisely what value ends up in a produced
\code{Kmer}'s \code{r} field, since it is easy to get this off by one and it is exercised by the
reference test suite. For the ($0$-based) $i$-th $k$-mer of a sequence (spanning nucleotide
positions $[i, i{+}k)$):
\[
  \texttt{Kmer.r} = i + k - 1
\]
i.e., \textbf{the index of the $k$-mer's last nucleotide}, not the exclusive upper bound of its
span (despite a doc-comment in the reference source claiming the latter). This is verified by the
reference test suite itself: for $k=4$ and a 5-character sequence, the two produced $k$-mers (at
$i=0,1$) have \code{r} $=3,4$ respectively, matching $i+k-1$.

\section{The rolling-hash recurrence}
\label{sec:fracminhash-hash}

Let $s$ be the input, $k$ the $k$-mer length, and define, for a byte $c$, $\mathrm{fw}(c) =
\text{\code{LUT\_fw}}[c]$ and $\mathrm{rc}(c) = \text{\code{LUT\_rc}}[c]$. The threshold is
\[
  H_{\mathrm{thr}} = \begin{cases} \lfloor r \cdot (2^{64}-1) \rfloor & r < 1 \\ 2^{64}-1 & r = 1
  \end{cases}
\]
computed once via a floating-point multiply-then-truncate (this is the one place floating-point
rounding can matter at all --- for $r=1$ it is special-cased to the exact maximum, avoiding any
rounding).

\paragraph{Initialization (first window, $i=0$).} Seed $h_{fw}=h_{rc}=0$, then fold in each of the
first $k$ bases with a position-dependent rotation:
\[
  h_{fw} \mathrel{\oplus}= \mathrm{rotl}\bigl(\mathrm{fw}(s[j]),\; k-j-1\bigr), \qquad
  h_{rc} \mathrel{\oplus}= \mathrm{rotl}\bigl(\mathrm{rc}(s[j]),\; j\bigr)
  \qquad\text{for } j = 0,\dots,k-1.
\]
(\(\mathrm{rotl}(x,n)\) / \(\mathrm{rotr}(x,n)\) denote 64-bit bitwise left/right rotation.) Note
the two hashes rotate in \emph{opposite} directions and by \emph{mirrored} amounts ($k{-}j{-}1$
vs.\ $j$) --- this is precisely what makes $h_{rc}$ equal to the forward hash of the
reverse-complement string: reversing the string reverses which position gets which rotation
amount, and complementing each base is handled already by using \code{LUT\_rc} instead of
\code{LUT\_fw}.

\paragraph{Emitting a $k$-mer.} At every window position (including the very first, before any
roll), combine the two sides and test the threshold:
\[
  h = h_{fw} \oplus h_{rc}, \qquad \text{keep iff } h \le H_{\mathrm{thr}}.
\]
If kept, its strand bit is $\;\texttt{strand} = (h_{fw} > h_{rc})\;$ --- i.e., whichever side is
numerically larger is deemed to have ``won'' and determines orientation. (This is an arbitrary
but self-consistent convention: swap the comparison and every strand bit flips consistently, with
no effect on any similarity computation, since those only ever compare two \code{Kmer}s'/\code{
Hit}s' strand bits to each other, never to an absolute notion of ``forward''.)

\paragraph{Rolling to the next window ($i \to i+1$, incoming base $s[i{+}k]$, outgoing base
$s[i]$).}
\[
  h_{fw} \leftarrow \mathrm{rotl}(h_{fw}, 1) \;\oplus\; \mathrm{rotl}(\mathrm{fw}(s[i]), k)
  \;\oplus\; \mathrm{fw}(s[i{+}k])
\]
\[
  h_{rc} \leftarrow \mathrm{rotr}(h_{rc}, 1) \;\oplus\; \mathrm{rotr}(\mathrm{rc}(s[i]), 1)
  \;\oplus\; \mathrm{rotl}(\mathrm{rc}(s[i{+}k]), k-1)
\]
Each update: (1) rotates the whole accumulated hash by one position in the direction that side
rolls, (2) removes the outgoing base's contribution (itself rotated to the position it now needs
to cancel from), (3) adds the incoming base's contribution at its new position. This is the
standard structure of an XOR/rotate rolling hash: rotation distributes each base's contribution
across all 64 bits over the $k$-window it participates in, and since $\mathrm{rotl}$/$\mathrm{rotr}$
are self-inverse under repeated application by the same total, the outgoing base's contribution
cancels \emph{exactly} (XOR with itself) once it has rolled all the way through.

\begin{algobox}{title={\code{sketch}(s, k, h\_frac)}}
\begin{verbatim}
if len(s) < k: return []
H_thr := (h_frac < 1) ? floor(h_frac * (2^64 - 1)) : 2^64 - 1
h_fw, h_rc := 0, 0
for j in 0..k-1:
    h_fw ^= rotl(LUT_fw[s[j]], k-j-1)
    h_rc ^= rotl(LUT_rc[s[j]], j)
r := k - 1              # position of the k-mer's last base, 0-based
result := []
loop:
    h := h_rc ^ h_fw
    if h <= H_thr:
        strand := (h_fw > h_rc)
        result.append(Kmer{ r: r, h: h, strand: strand })
    if r + 1 >= len(s): break
    out_c, in_c := s[r + 1 - k], s[r + 1]
    h_fw := rotl(h_fw, 1) ^ rotl(LUT_fw[out_c], k) ^ LUT_fw[in_c]
    h_rc := rotr(h_rc, 1) ^ rotr(LUT_rc[out_c], 1) ^ rotl(LUT_rc[in_c], k-1)
    r := r + 1
return result
\end{verbatim}
\end{algobox}

\begin{notebox}
The loop above is written with the loop variable \code{r} directly holding the emitted
\code{Kmer.r} value at every iteration (equivalent to, but clearer than, the reference source's
own \code{for(r=0;...) ... push(Kmer(r-1,...))} formulation, which tracks ``one past'' internally
and subtracts 1 only at emission time). Both are exactly equivalent; verify against
\S\ref{sec:sketch-r-semantics} if re-deriving independently.
\end{notebox}

\section{Reference C++ implementation}

\begin{lstlisting}[style=cppstyle,caption={\code{FracMinHash::sketch}, \texttt{src/sketch.h} (verbatim)}]
const sketch_t sketch(const std::string& s) const {
    sketch_t kmers;
    kmers.reserve((rpos_t)(1.1 * (double)s.size() * hFrac));
    if (rpos_t(s.size()) < k) return kmers;

    hash_t h_fw = 0, h_rc = 0;
    hash_t hThres = hFrac < 1.0
        ? hash_t(hFrac * double(std::numeric_limits<hash_t>::max()))
        : std::numeric_limits<hash_t>::max();
    rpos_t r;

    for (r=0; r<k; r++) {
        h_fw ^= std::rotl(LUT_fw[ size_t(s[r]) ], k-r-1);
        h_rc ^= std::rotl(LUT_rc[ size_t(s[r]) ], r);
    }

    while(true) {
        auto h = h_rc ^ h_fw;
        if (h <= hThres) {
            const bool strand = h_fw > h_rc;
            kmers.push_back(Kmer(r-1, h, strand));
        }
        if (r >= rpos_t(s.size())) break;
        h_fw = std::rotl(h_fw, 1) ^ std::rotl(LUT_fw[ size_t(s[r-k]) ], k) ^ LUT_fw[ size_t(s[r]) ];
        h_rc = std::rotr(h_rc, 1) ^ std::rotr(LUT_rc[ size_t(s[r-k]) ], 1)
             ^ std::rotl(LUT_rc[ size_t(s[r]) ], k-1);
        ++r;
    }
    return kmers;
}
\end{lstlisting}

\section{Rust re-implementation}

\begin{lstlisting}[style=rststyle,caption={\code{FracMinHash::sketch}, Rust port (\texttt{src/sketch.rs})}]
pub fn sketch(&self, s: &[u8], counters: &mut Counters) -> SketchT {
    let k = self.k;
    let mut kmers: SketchT =
        Vec::with_capacity((1.1 * s.len() as f64 * self.h_frac).max(0.0) as usize);
    if (s.len() as RPos) < k { return kmers; }

    let h_thres: Hash = if self.h_frac < 1.0 {
        (self.h_frac * u64::MAX as f64) as u64
    } else { u64::MAX };

    let mut h_fw: Hash = 0;
    let mut h_rc: Hash = 0;
    let mut r: RPos = 0;
    while r < k {
        let c = s[r as usize] as usize;
        h_fw ^= self.lut_fw[c].rotate_left((k - r - 1) as u32);
        h_rc ^= self.lut_rc[c].rotate_left(r as u32);
        r += 1;
    }

    loop {
        let h = h_rc ^ h_fw;
        if h <= h_thres {
            let strand = h_fw > h_rc;
            kmers.push(Kmer::new(r - 1, h, strand));
        }
        if r >= s.len() as RPos { break; }
        let out_c = s[(r - k) as usize] as usize;
        let in_c = s[r as usize] as usize;
        h_fw = h_fw.rotate_left(1) ^ self.lut_fw[out_c].rotate_left(k as u32) ^ self.lut_fw[in_c];
        h_rc = h_rc.rotate_right(1)
            ^ self.lut_rc[out_c].rotate_right(1)
            ^ self.lut_rc[in_c].rotate_left((k - 1) as u32);
        r += 1;
    }
    kmers
}
\end{lstlisting}

Note the Rust version zero-initializes both LUTs before filling the 8 real entries (the fix from
the bug box above); everything else is a direct, bit-for-bit translation --- \code{rotate\_left}/
\code{rotate\_right} on a \code{u64} are exactly \code{std::rotl}/\code{std::rotr} on
\code{uint64\_t}, including the fact that both are defined for \emph{any} shift amount (taken mod
64), so no extra bounds-masking is needed even though $k$ can exceed typical small shift-amount
assumptions.

\section{Non-ACGT input}
\label{sec:sketch-nonacgt}

Real sequencing data occasionally contains \texttt{N} (unknown base) or, rarely, IUPAC ambiguity
codes. As established above, the reference implementation's behavior on these is undefined; the
recommended, safe, deterministic choice --- adopted by the Rust port --- is that such a byte
contributes a hash of exactly $0$ on both the forward and reverse-complement side. This is a
reasonable default because it is simple, deterministic, and does not privilege any particular
real base; a from-scratch implementation is free to choose a different deterministic convention
(e.g.\ skip $k$-mers containing any non-ACGT byte entirely, which several other tools do) as long
as it is applied \emph{identically} when sketching the reference and the reads, since any
asymmetry there would silently break matching.

\section{Complexity}

Sketching a sequence of length $n$ takes $O(n)$ time (one $O(1)$-amortized rolling-hash step per
position, after an $O(k)$ initialization) and produces, in expectation, $O(rn)$ output $k$-mers,
using $O(rn)$ space for the output (input is streamed, not buffered beyond what the caller already
holds).

\section{Worked example}

For $k=4$, $r=1.0$ (keep everything) and input \texttt{"ACGGT"} (length 5, so 2 $k$-mers), the two
produced \code{Kmer}s have \code{r} $= 3, 4$ (matching $i+k-1$ for $i=0,1$). Feeding the
reverse complement of the same underlying sequence, \texttt{"ACCGT"}, produces the \emph{same
multiset of hashes} in reversed position order and with each \code{strand} bit flipped relative to
the original --- this symmetry (verified as an automated test in both the C++ and Rust
implementations) is the single most important sanity check to run against any from-scratch
reimplementation before trusting it: sketch a string and its reverse complement, reverse one
result, and confirm the hash sequences match exactly.
